\documentclass[11pt]{article}

\usepackage[margin=1in]{geometry}
\usepackage{amsmath,amssymb,amsthm}
\usepackage{mathtools}
\usepackage{enumitem}
\usepackage{hyperref}

\theoremstyle{plain}
\newtheorem{theorem}{Theorem}
\newtheorem{lemma}[theorem]{Lemma}
\newtheorem{corollary}[theorem]{Corollary}
\newtheorem{definition}[theorem]{Definition}
\newtheorem{remark}[theorem]{Remark}

% ---------------------------------------------------------------------
% Macros
% ---------------------------------------------------------------------
\newcommand{\Ssem}[1]{\mathsf{S}[\![#1]\!]}
\newcommand{\Lsem}[1]{\mathsf{L}[\![#1]\!]}
\newcommand{\Vsem}[1]{\mathsf{V}[\![#1]\!]}

\newcommand{\Wm}[1]{\widehat{W}_{#1}}
\newcommand{\Rm}[1]{\widehat{R}_{#1}}
\newcommand{\Sm}[1]{\widehat{S}_{#1}}

\newcommand{\eps}{\epsilon}
\newcommand{\dg}[1]{^{\dagger #1}}
\newcommand{\op}{\mathit{op}}
\newcommand{\then}{\mathrel{\blacktriangleright}}
\newcommand{\loc}[1]{\langle #1 \rangle} % local construct, decorated separately
\newcommand{\hoist}{\mathsf{hoist}}
\newcommand{\tell}{\mathsf{tell}}
\newcommand{\censor}{\mathsf{censor}_0}
\newcommand{\collect}{\triangle}
\newcommand{\collectX}[1]{\mathsf{collect}_{#1}}
\newcommand{\collapse}{\mathsf{collapse}}
\newcommand{\bump}{\mathsf{bump}}
\newcommand{\Rx}[1]{\widehat{\mathcal R}_{#1}}
\newcommand{\lC}{\lambda_C}

\title{An Alternative Denotational Semantics for $\lambda_C$:\\
Swapping the Free Monad and the Writer Monad Transformer}
\author{Companion to ``Handling the Selection Monad''\\[2pt]
\small (after Gordon Plotkin and Ningning Xie)}
\date{}

\begin{document}
\maketitle

\begin{abstract}
Section~5 of Plotkin and Xie's \emph{Handling the Selection Monad} gives $\lambda_C$ a
denotational semantics in which the auxiliary monad $W_\eps(X) = F_\eps(R\times X)$ arises
by applying the \emph{writer monad transformer} to the free monad $F_\eps$: effect trees are
built first, and a single accumulated loss is attached at each leaf. This note develops the
mirror-image construction, in which the \emph{free monad transformer} is applied to the
\emph{writer monad}: $\widehat W_\eps(X) = \mathsf{FreeT}_\eps(\mathsf{Writer}_R)(X)$, so
that a loss contribution is attached at \emph{every} node of an effect tree rather than only
at its leaves. We re-derive the monad structure, the selection-monad construction
$\widehat S_\eps$, the semantics of expressions, and the semantics of handlers for this
alternative choice, and discuss the one substantive semantic difference it introduces: loss
no longer commutes with operation calls. The construction is exactly the one realised by the
reference implementation \texttt{Sel r e a = FreeT e (Writer r)}, which we use throughout as
a sanity check on the mathematics.
\end{abstract}

\section{Introduction}

Section 5.1 of the paper introduces the family of augmented selection monads
$S_\eps(X) = (X\to R_\eps)\to W_\eps(X)$, where
\[
  W_\eps(X) = F_\eps(R\times X), \qquad R_\eps = F_\eps(R),
\]
and $F_\eps$ is the free monad for the effect signature at level $\eps$. The paper observes
that $W_\eps$ is ``the commutative combination [Hyland et al.\ 2006] of $F_\eps$ and the
writer monad $R\times -$'', and that this choice makes the loss operation \emph{commute}
with every other operation --- semantically, a loss can be freely floated past an operation
call, and this is exactly what the operational semantics does (Theorem~5.4).

Categorically, $F_\eps(R\times X)$ is what one obtains by applying the \emph{writer monad
transformer} $\mathsf{WriterT}_R$ to the base monad $F_\eps$: a $\mathsf{WriterT}_R\, M$
computation is an $M$-computation of a value paired with a monoid element, i.e.\ $M(R\times
X)$, and here $M = F_\eps$. This note asks the dual question: what if instead we apply the
\emph{free monad transformer} $\mathsf{FreeT}_\eps$ to the base monad
$\mathsf{Writer}_R = R\times -$? Concretely we replace
\[
  W_\eps(X) = F_\eps(R \times X)
  \qquad\text{by}\qquad
  \widehat W_\eps(X) = \mathsf{FreeT}_\eps(\mathsf{Writer}_R)(X),
\]
i.e.\ we thread a writer layer through \emph{every} constructor of the effect tree, rather
than combining all the writer output into a single pair at the leaves. This is precisely the
monad implemented in the accompanying Haskell development (file \texttt{Correct.hs}):
\begin{quote}
\texttt{newtype Sel r e a = Sel \{ \_runSel :: (a -> FreeT e (Writer r) ()) -> FreeT e (Writer r) a \}}
\end{quote}
in which the carrier of $W_\eps$ is literally \texttt{FreeT e (Writer r) a}. We use that
implementation throughout as ground truth for the constructions below, presenting them in
the mathematical style of the original paper's Section~5 so that the two developments can be
read side by side.

Everything in Section~5.1 that concerns the semantics of \emph{types} (Fig.~8) is orthogonal
to this swap and carries over unchanged; we recall it briefly in Section~2. Section~3
develops the new monad $\widehat W_\eps$ and the resulting selection monad $\widehat S_\eps$.
Section~4 redoes the semantics of expressions (Fig.~9). Section~5 redoes the semantics of
handlers. Section~6 discusses the one genuine semantic difference (loss no longer commutes
with effects) and what it implies for soundness and adequacy.

\section{Recap: Semantics of Types}

We keep the syntax of $\lambda_C$, its typing rules, and its operational semantics exactly as
in Sections 3--4 of the original paper: types $\sigma,\tau$, effect contexts $\eps$ (multisets
of effect labels), a signature $\Sigma$ of operations $\op : \mathit{out}\xrightarrow{\ell}
\mathit{in}$ for $\ell\in\eps$, expressions $e$, loss-continuation expressions $g$, and
handlers $h$ with clauses $\op_i\mapsto\dots$ and $\mathit{return}\mapsto\dots$.

The semantics of types, $\Sm{\sigma}$, is exactly Figure~8 of the original paper:
\[
  \Ssem{b} = [\![b]\!], \qquad
  \Ssem{(\sigma_1,\dots,\sigma_n)} = \Ssem{\sigma_1}\times\cdots\times\Ssem{\sigma_n}, \qquad
  \Ssem{\sigma\to\tau\,!\,\eps} = \Ssem{\sigma}\to \Sm\eps(\Ssem{\tau}),
\]
where $\Sm\eps$ is now the family we redefine in Section~3. Only the auxiliary monad
changes; the interpretation of type formers does not.

\section{The Alternative Monad}

\subsection{The carrier $\widehat W_\eps$}

Recall an $\eps$-algebra is a set $X$ with maps
$\varphi_{\ell,\op,i} : \Ssem{\mathit{out}}\times X^{\Ssem{\mathit{in}}}\to X$ for each
$\ell\in\eps$, $\op:\mathit{out}\xrightarrow{\ell}\mathit{in}$, $0<i\le\eps(\ell)$, and that
$F_\eps(X)$ is the free such algebra on $X$. We now dispense with a separate ``plain'' free
monad: because the writer layer is woven through every node rather than paired only at
leaves, one operator suffices where the original development needed two ($F_\eps$ and
$W_\eps = F_\eps(R\times-)$).

\begin{definition}[The monad $\widehat W_\eps$]
For a set $X$, let $\widehat W_\eps(X)$ be the least set $Y$ such that
\[
  Y \;=\; R \times \left(X \;+\; \sum_{\substack{\ell\in\eps,\ \op:\mathit{out}\xrightarrow{\ell}\mathit{in}\\ 0<i\le\eps(\ell)}}
        \Ssem{\mathit{out}} \times Y^{\Ssem{\mathit{in}}}\right).
\]
We write an element as a pair $(r,n)$ with $r\in R$ (the loss recorded \emph{at this node})
and $n$ either a leaf $\iota_1(x)$ for $x\in X$, or an internal node
$\iota_2\big((\ell,\op,i),(o,\kappa)\big)$ decorated exactly as in $F_\eps$.
\end{definition}

This is exactly the least-fixpoint reading of the Haskell type
\texttt{FreeT e (Writer r) a}, whose definition
\begin{quote}
\texttt{data FreeF f a b = Pure a | Free (f b)}\\
\texttt{newtype FreeT f m a = FreeT \{ runFreeT :: m (FreeF f a (FreeT f m a)) \}}
\end{quote}
instantiated at $m = \mathsf{Writer}_R$ unrolls to $\mathsf{Writer}_R(X + F(\widehat
W_\eps(X))) = R\times(X + F(\widehat W_\eps(X)))$, matching the equation above with $F$ the
signature functor of $\eps$.

Compare with the original $W_\eps(X) = F_\eps(R\times X)$, whose defining equation is
\[
  Y \;=\; \sum_{\ell,\op,i} \Ssem{\mathit{out}}\times Y^{\Ssem{\mathit{in}}} \;+\; (R\times X):
\]
there $R$ occurs exactly once, paired with the eventual result, at each leaf; here $R$ occurs
at \emph{every} node, including every internal one. As before, the two families $\widehat
W_\eps$ and (the semantics $\widehat S_\eps$ defined below) are mutually recursive, with
effect levels strictly decreasing on one side and non-increasing on the other exactly as
described in the original Section~5.1; that circularity argument is untouched by the swap, so
we do not repeat it.

We write $\Rm\eps := \widehat W_\eps(1)$ for the type of \emph{loss continuations}, exactly
mirroring $R_\eps := W_\eps(1) \cong F_\eps(R)$ in the original -- except that here, unlike in
the original, $\widehat W_\eps(1)$ is \emph{not} isomorphic to a simple type of accumulated
losses: it retains the full tree shape (an unhandled operation can still be pending at any
depth), it merely carries a unit value at each leaf.

\subsection{Free layered $\eps$-algebras}

The original development makes $W_\eps$ a monad by observing that it is the free \emph{action
$\eps$-algebra}: an $\eps$-algebra $(Y,\psi)$ together with an additive action $\cdot : R\times
Y\to Y$ ($0\cdot y = y$, $r\cdot(s\cdot y)=(r+s)\cdot y$) that \emph{commutes} with every
$\psi_{\ell,\op,i}$. That commutation is exactly what makes the combination ``commutative''
in the sense of Hyland, Plotkin and Power, and exactly what the swap gives up.

\begin{definition}[Layered $\eps$-algebra]
A \emph{layered $\eps$-algebra} is a set $Y$ equipped with an $\eps$-algebra structure $\psi$
and an additive action $\cdot: R\times Y \to Y$, with \emph{no} requirement that $\cdot$
commutes with $\psi$.
\end{definition}

\begin{definition}[Free layered $\eps$-algebra]
$\widehat W_\eps(X)$ is the free layered $\eps$-algebra on $X$: its own structure is
\[
  \widehat\varphi_{\ell,\op,i}(o,\kappa) =_{\mathrm{def}} \big(0,\, \iota_2((\ell,\op,i),(o,\kappa))\big),
  \qquad
  r\cdot(r_0,n) =_{\mathrm{def}} (r+r_0,\,n),
\]
with unit $\eta^{\widehat W_\eps}(x) = (0,\iota_1(x))$; and for any layered $\eps$-algebra
$(Y,\psi,\cdot)$ and function $f:X\to Y$, the unique homomorphic extension
$f\dg{\widehat W_\eps}: \widehat W_\eps(X)\to Y$ is given by
\[
  f\dg{\widehat W_\eps}(r_0,\iota_1(x)) = r_0\cdot f(x),
  \qquad
  f\dg{\widehat W_\eps}\big(r_0,\iota_2((\ell,\op,i),(o,\kappa))\big)
     = r_0\cdot \psi_{\ell,\op,i}\big(o,\ f\dg{\widehat W_\eps}\circ\kappa\big).
\]
\end{definition}

Both clauses are well defined by structural recursion since $\widehat W_\eps(X)$ is a least
(hence well-founded, finitely-branching-and-depth) set. One checks directly that
$f\dg{\widehat W_\eps}\circ\eta^{\widehat W_\eps} = f$ and that $f\dg{\widehat W_\eps}$ is a
homomorphism of layered $\eps$-algebras, using only the additive action law of $Y$ -- no
commutation between $\psi$ and $\cdot$ is needed anywhere in the argument, which is precisely
why $\widehat\varphi$ and the action need not commute for the recursion to be well founded.

\begin{remark}
Instantiating $(Y,\psi,\cdot) := (\widehat W_\eps(Z),\widehat\varphi,\cdot)$ for a function
$f : X\to \widehat W_\eps(Z)$ turns the homomorphic-extension formula above into the
\emph{Kleisli extension} (bind) of $\widehat W_\eps$: writing $\mathsf{let}_{\widehat
W_\eps}\, x\in X\ \mathsf{be}\ u\ \mathsf{in}\ f(x) := f\dg{\widehat W_\eps}(u)$,
\[
  \mathsf{let}_{\widehat W_\eps}\ x\in X\ \mathsf{be}\ (r_0,\iota_1(x'))\ \mathsf{in}\ f(x)
     = r_0\cdot f(x'),
\]
\[
  \mathsf{let}_{\widehat W_\eps}\ x\in X\ \mathsf{be}\ (r_0,\iota_2((\ell,\op,i),(o,\kappa)))\ \mathsf{in}\ f(x)
     = \Big(r_0,\ \iota_2\big((\ell,\op,i),\ (o,\ \lambda a.\ \mathsf{let}_{\widehat W_\eps}\ x\in X\ \mathsf{be}\ \kappa(a)\ \mathsf{in}\ f(x))\big)\Big).
\]
That is: at a leaf, the node's own loss $r_0$ is \emph{added} to whatever the continuation
produces (via the action); at an internal node, $r_0$ is left exactly where it is and only the
children are rebound. This is exactly the standard bind for a monad transformer stacked over
$\mathsf{Writer}_R$, and matches the (derivable) \texttt{Monad} instance underlying
\texttt{FreeT e (Writer r)} in \texttt{Correct.hs}.
\end{remark}

\subsection{Loss, commutativity, and what changes}

The reason the original $W_\eps$ needs the stronger notion of \emph{action} $\eps$-algebra
(with $\psi$ and $\cdot$ commuting) is that there loss is combined into a single value at the
very end; an operation occurring ``after'' a loss and one occurring ``before'' it produce
indistinguishable trees, which is exactly the content of commutativity. Here, because a loss
is recorded at the specific node at which it occurs, the order of a loss relative to a
neighbouring operation is part of the data of $\widehat W_\eps(X)$: $\widehat\varphi$ applied
first and then $\tell$ (defined below) applied to the result is a \emph{different} element
from $\tell$ applied first and then $\widehat\varphi$. Section~6 returns to the operational
consequences of this; here we simply record the two combinators we will need:
\[
  \tell(r)(r_0,n) =_{\mathrm{def}} (r+r_0,\,n) \qquad(\text{i.e. } \tell(r) = r\cdot(-)),
\]
\[
  \censor(r_0,n) =_{\mathrm{def}}
  \begin{cases}
    (0,\iota_1(x)) & n=\iota_1(x)\\
    (0,\iota_2((\ell,\op,i),(o,\lambda a.\,\censor(\kappa(a))))) & n=\iota_2((\ell,\op,i),(o,\kappa)),
  \end{cases}
\]
the latter recursively zeroing every recorded loss throughout the tree while leaving its shape
and leaves untouched (matching Haskell's \texttt{censor (const mempty)}), and
\[
  \collect(r_0,n) =_{\mathrm{def}}
  \begin{cases}
    (0,\iota_1(r_0)) & n = \iota_1(())\\
    \tell(r_0)\big(0,\iota_2((\ell,\op,i),(o,\lambda a.\,\collect(\kappa(a))))\big) & n = \iota_2((\ell,\op,i),(o,\kappa)),
  \end{cases}
\]
a map $\widehat W_\eps(1)\to \widehat W_\eps(R)$ that recursively \emph{sums} all losses along
each root-to-leaf path into that leaf's now $R$-valued payload, while zeroing every node's own
recorded loss along the way (matching Haskell's
\texttt{triangle = censor (const mempty) . fmap snd . listen}). We also need the evident
generalisation to arbitrary payload type $X$, matching Haskell's \texttt{listen}, which pairs
each leaf value with the total accumulated loss on its path rather than discarding it:
\[
  \collectX X(r_0,\iota_1(x)) = (0,\iota_1((x,r_0))),
  \qquad
  \collectX X\big(r_0,\iota_2((\ell,\op,i),(o,\kappa))\big)
     = \Big(0,\ \iota_2\big((\ell,\op,i),\ (o,\ \lambda a.\ \bump_{r_0}(\collectX X(\kappa(a))))\big)\Big),
\]
where $\bump_r : \widehat W_\eps(X\times R)\to\widehat W_\eps(X\times R)$ adds $r$ into the
recorded total at \emph{every} leaf beneath the current node (needed because, unlike $\tell$,
$\collectX X$ must push a node's loss down into leaves that may be arbitrarily deep, not just
combine it at the top):
\[
  \bump_r\big(0,\iota_1((x,s))\big) = \big(0,\iota_1((x,r+s))\big),
  \qquad
  \bump_r\big(0,\iota_2((\ell,\op,i),(o,\kappa))\big) = \big(0,\iota_2((\ell,\op,i),(o,\lambda a.\,\bump_r(\kappa(a))))\big)
\]
(the outer loss is always $0$ on the nose in the argument of $\bump_r$, since $\collectX X$
zeroes it just before recursing). One checks $\bump_0 = \mathrm{id}$ and
$\bump_r\circ\bump_s = \bump_{r+s}$ directly from these clauses. Finally
$\collect = \Wm\eps(\pi_2)\circ\collectX 1 : \widehat W_\eps(1)\to\widehat W_\eps(R)$, i.e.\
$\collect$ is $\collectX 1$ followed by forgetting the (unit-typed) leaf value and keeping only
the accumulated loss -- matching \texttt{triangle}'s definition as \texttt{censor (const
mempty) . fmap snd . listen} exactly, clause by clause.

\subsection{The selection monad $\widehat S_\eps$}

With $\widehat W_\eps$ and $\Rm\eps=\widehat W_\eps(1)$ in hand, the construction of the
selection monad itself is \emph{generic} in the choice of $\widehat W_\eps$: it is exactly
the scheme of the original paper's equation~(5), reinstantiated. Define
\[
  \widehat S_\eps(X) = (X\to \Rm\eps)\to \widehat W_\eps(X), \qquad
  \eta^{\widehat S_\eps}(x) = \lambda\gamma.\ \eta^{\widehat W_\eps}(x).
\]
$\Rm\eps$ is itself a layered $\eps$-algebra (being $\widehat W_\eps(1)$), with action
$r\cdot u = \tell(r)(u)$. For $F\in\widehat S_\eps(Y)$ and a loss function
$\gamma:Y\to\Rm\eps$, define, exactly as in the original,
\[
  \widehat{\mathcal R}_\eps(F\mid\gamma) =_{\mathrm{def}} \gamma\dg{\widehat W_\eps}(F(\gamma)) \ \in\ \Rm\eps.
\]
Then the Kleisli extension of $f:X\to\widehat S_\eps(Y)$ is, verbatim,
\[
  f\dg{\widehat S_\eps}(F) = \lambda\gamma\in Y\to\Rm\eps.\
     \mathsf{let}_{\widehat W_\eps}\ x\in X\ \mathsf{be}\ F\big(\lambda x\in X.\ \widehat{\mathcal R}_\eps(fx\mid\gamma)\big)\ \mathsf{in}\ fx\,\gamma,
\]
and $\widehat S_\eps(X)$ is itself an $\eps$-algebra via
\[
  \widehat\varphi^X_{\ell,\op,i}(o,f)(\gamma) = \widehat\varphi_{\ell,\op,i}\big(o,\ \lambda x\in\mathit{in}.\ f(x)(\gamma)\big).
\]
Two simplifications are worth flagging, both consequences of no longer needing to
pre-multiply values by $R$ before feeding them to $F_\eps$: the $\varphi$ used in the last
equation is the one-argument $\widehat\varphi$ of $\widehat W_\eps(X)$ directly, rather than
$\varphi^{R\times X}$; and nothing elsewhere in this subsection ever mentions $R\times -$
explicitly. All bookkeeping of losses is delegated to the fact that $\widehat W_\eps$ itself
already carries $R$ at every node.

\section{Semantics of Expressions}

Most of Figure~9 is stated purely in terms of the abstract interface of $S_\eps$ ($\eta$,
Kleisli extension, and $\varphi^X$) and therefore carries over to $\widehat S_\eps$
\emph{unchanged in form}: the clauses for variables, constants, tuples, projection,
abstraction, application, and operation calls
\[
  \Ssem{\op(e)}(\rho) = \mathsf{let}_{\widehat S_\eps}\ a\in\Ssem{\mathit{out}}\ \mathsf{be}\ \Ssem{e}(\rho)
     \ \mathsf{in}\ \widehat\varphi^X_{\ell,\op,\eps(\ell)}\big(a,(\eta^{\widehat S_\eps})^{\Ssem{\mathit{in}}}\big)
\]
are exactly as in the original, read with $\widehat S_\eps,\widehat W_\eps,\widehat
\varphi$ in place of $S_\eps,W_\eps,\varphi$. The four constructs that manipulate the
$F_\eps(R\times X)$ leaf shape directly -- \texttt{loss}, $(\then)$, the local construct
$\loc{e}g^\eps$, and \texttt{reset} -- need to be redone against the new shape. We also
restate the auxiliary loss-function semantics $\mathcal L$.

\paragraph{\texttt{loss(e)}.} Originally this pattern-matched the leaf pair of
$\Ssem{e}(\rho)(\gamma)\in F_\eps(R\times R)$ and added the returned value into the recorded
loss. Here, since a bare $\tell$ is already an $\widehat S_\eps$-level operation, no pattern
matching is needed at all:
\[
  \Ssem{\mathtt{loss}(e)}(\rho) = \lambda\gamma.\
    \mathsf{let}_{\widehat W_\eps}\ a\in R\ \mathsf{be}\ \Ssem{e}(\rho)(\gamma)\ \mathsf{in}\
    \tell(a)\big(\eta^{\widehat W_\eps}(())\big).
\]

\paragraph{Loss-function semantics $\mathcal L$.} For $\Gamma,x\!:\!\sigma\vdash e:\mathtt{loss}\,!\,\eps$,
define $\Lsem{\lambda^\eps x\!:\!\sigma.\,e} : \Sm\Gamma \to \Sm\sigma \to \Rm\eps$ by
\[
  \Lsem{\lambda^\eps x\!:\!\sigma.\,e}(\rho) = \lambda a\in\Ssem\sigma.\
     \big(\lambda(r_0,n).\ (0,n)\big)\Big[\, v \mapsto v \,\Big]\big(\Ssem{e}(\rho[a/x])(\gamma_0)\big)
\]
which, spelled out without the shorthand, is exactly the recursive map
\[
  \Lsem{\lambda^\eps x\!:\!\sigma.\,e}(\rho)(a) =
  \begin{cases}
    (v,\iota_1(())) & \Ssem{e}(\rho[a/x])(\gamma_0) = (r_0,\iota_1(v)),\ v\in R\\
    \big(0,\iota_2((\ell,\op,i),(o,\lambda w.\,\Lsem{\dots}(\rho)(a)\!\downharpoonright_{\kappa(w)}))\big) & \text{on internal nodes},
  \end{cases}
\]
i.e.\ the analogue of the original's ``discard the accumulated loss $r_1$, keep only the
returned value $r_2$'', recursing structurally through any still-pending operation nodes and
zeroing their recorded losses on the way. ($\gamma_0 = \lambda r{\in}R.\,\eta^{\widehat
W_\eps}(())$ is the zero loss continuation, exactly as before.) We call this operation
$\mathsf{collapse}: \widehat W_\eps(R)\to\widehat W_\eps(1)$ below and simply write
$\Lsem{\lambda^\eps x\!:\!\sigma.\,e}(\rho)(a) = \mathsf{collapse}\big(\Ssem e(\rho[a/x])(\gamma_0)\big)$.

\paragraph{$e_1\then\lambda^{\eps_1} x\!:\!\sigma_1.\,e_2$.} This builds the loss recorded for
one step of a loss continuation: run $e_1$ under the hypothetical continuation
$\Lsem{\lambda x.\,e_2}(\rho)$, obtaining a value $a$ together with the loss $r_1$ this
incurred; then run $e_2$ at $a$ under the zero continuation to get its own bookkeeping loss
$r_2$ and returned increment $r_3$; the result records loss $r_2$ and returns $r_1+r_3$. Using
$\mathsf{collect}_X$ from Section 3.3 to extract $(a,r_1)$ pairs from a tree of pending
operations,
\[
  \Ssem{e_1\then\lambda^{\eps_1}x\!:\!\sigma_1.\,e_2}(\rho) = \lambda\gamma.\
    \mathsf{let}_{\widehat W_{\eps_1}}\ (a,r_1)\in\Ssem{\sigma_1}\times R\ \mathsf{be}\
      \mathsf{collect}_{\Ssem{\sigma_1}}\big(\Ssem{e_1}(\rho)(\Lsem{\lambda x.\,e_2}(\rho))\big)\ \mathsf{in}
\]
\[
  \hphantom{\Ssem{e_1\then\lambda^{\eps_1}x\!:\!\sigma_1.\,e_2}(\rho) = }
    \mathsf{let}_{\widehat W_{\eps_1}}\ (r_3,r_2)\in R\times R\ \mathsf{be}\
      \mathsf{collect}_R\big(\Ssem{e_2}(\rho[a/x])(\gamma_0)\big)\ \mathsf{in}\
      \tell(r_2)\big(\eta^{\widehat W_{\eps_1}}(r_1+r_3)\big).
\]

\paragraph{$\loc{e}g^{\eps_1}$.} Exactly as in the original, this clause only substitutes a
different loss continuation and is oblivious to the concrete shape of $\widehat W_\eps$, so it
transfers verbatim:
\[
  \Ssem{\loc{e}g^{\eps_1}}(\rho) = \lambda\gamma.\ \Ssem{e}(\rho)\big(\Lsem{g}(\rho)\big).
\]

\paragraph{\texttt{reset}\,$e$.} Originally: run $e$ under $\gamma$, discard the recorded loss
entirely, keep the value. Here that is exactly $\censor$ from Section~3.3:
\[
  \Ssem{\mathtt{reset}\ e}(\rho) = \lambda\gamma.\ \censor\big(\Ssem{e}(\rho)(\gamma)\big).
\]
This matches \texttt{Correct.hs} exactly: \texttt{reset p = Sel (\textbackslash g ->
censor (const mempty) (runSel g p))}.

\section{Semantics of Handlers}

We follow the same staged construction as the original (Section 5.3): fix $\rho\in\Sm\Gamma$
and $\gamma\in\Sm{\sigma'}\to\Rm\eps$, and build a layered $\eps\ell$-algebra structure on
\[
  A = \widehat W_\eps(\Sm{\sigma'})^{\Sm{\mathit{par}}}.
\]
Because the writer layer is now carried automatically by $\widehat W_\eps$ itself, the
construction needs no explicit multiplication by $R$ anywhere -- this is the one place where
the swap strictly \emph{simplifies} the original definitions. For the operations
\emph{not} handled by $h$ ($\ell_1\in\eps$, i.e.\ shallower than the newly-introduced
occurrence of $\ell$), pass-through:
\[
  \psi_{\ell_1,\op,i}(o,k) = \lambda p.\ \widehat\varphi_{\ell_1,\op,i}\big(o,\ \lambda a.\, k(a)(p)\big).
\]
For the handler's own clauses $\op_j\mapsto\lambda^\eps z\!:\!(\dots).\,e_j$, at the newly
introduced occurrence ($i = \eps(\ell)+1$):
\[
  \psi_{\ell,\op_j,i}(o,k) = \lambda p.\ \Ssem{e_j}\big(\rho[(p,o,l_1,k_1)/z]\big)(\gamma),
\]
where the delimited continuation is $k_1(p,a) = k(a)(p)$ exactly as before, and the choice
continuation -- now computed without any $\delta_\eps$-style leaf surgery, using $\collect$ in
place of it -- is
\[
  l_1(p,a) = \lambda\gamma_1.\ \collect\Big(\gamma\dg{\widehat W_\eps}\big(k(a)(p)(\gamma)\big)\Big)
  \quad \in \widehat S_\eps(R),
\]
i.e.\ it \emph{ignores} its nominal continuation $\gamma_1$ (matching Haskell's
\texttt{hoist}, which does the same), runs the delimited continuation under the ambient
$\gamma$, extends $\gamma$ homomorphically over the result to obtain a value of $\Rm\eps$, and
promotes that to an $R$-valued tree via $\collect$ so it can be used as a $\mathtt{loss}$-typed
value inside $e_j$. This is line for line
\begin{quote}
\texttt{hoist (triangle (go g y >>= g))}
\end{quote}
from \texttt{handle} in \texttt{Correct.hs}, once \texttt{go g} is unfolded as ``run the
handler with loss continuation \texttt{g}''.

The return clause needs no multiplication by an accumulated loss either -- again because
$\widehat W_\eps$ already threads it -- giving a strictly simpler $s$ than the original's
$s(r,a)=\lambda p.\,r\cdot(\dots)$:
\[
  s(a) = \lambda p.\ \Ssem{e_r}\big(\rho[(p,a)/z]\big)(\gamma) \ \in A,
\]
matching Haskell's \texttt{hret :: a -> Sel r e b}, which likewise takes no separate loss
argument. Finally, using the free-algebra property of $\widehat W_{\eps\ell}(\Sm\sigma)$ over
$A$ (Section 3.2, now instantiated directly at $X = \Sm\sigma$, with no $R\times$ pairing
needed to seed the extension):
\[
  \Ssem{h}(\rho)(p,G)(\gamma) = s\dg{\widehat W_{\eps\ell}}\Big(G\big(\lambda a\in\Sm\sigma.\
    \gamma\dg{\widehat W_\eps}(s(a)(p))\big)\Big)(p),
\]
where the inclusion $\Rm\eps\hookrightarrow\Rm{\eps\ell}$ (from $\eps\subseteq\eps\ell$) is
left implicit, exactly as sub-effect inclusions are left implicit throughout the original
paper. This is the mathematical reading of
\begin{quote}
\texttt{go g (runSel (\textbackslash x -> transFreeT RightEff (runSel g (hret h x) >>= g)) p)}
\end{quote}
where \texttt{transFreeT RightEff} implements the same inclusion $\Rm\eps\hookrightarrow
\Rm{\eps\ell}$, and \texttt{go g} implements $s\dg{\widehat W_{\eps\ell}}(-)(p)$.

\section{Discussion}

\paragraph{What the swap costs and buys.} Two things changed relative to the original
Section~5: first, $R$ no longer needs to be paired with the payload type before feeding it to
a free-monad construction, which removed a handful of $R\times(-)$ occurrences from the
$\widehat S_\eps$ scheme, the return clause $s$, and the handler extension; second, and more
substantively, the equations $\psi_{\ell,\op,i}$ and the action $\cdot$ on $\widehat W_\eps$ do
\emph{not} commute (Section~3.3). Concretely, evaluating \texttt{loss(e1); op(e2)} and
\texttt{op(e2); loss(e1)} (for independent $e_1,e_2$) now denote different elements of
$\widehat W_\eps$: the loss is attached to a different node in each case, whereas under the
original $F_\eps(R\times-)$ both denote the same tree with the same single leaf-loss. This is
the direct semantic mirror of the algebraic fact from Hyland, Plotkin and Power's
\emph{tensor}: the original $W_\eps$ is (isomorphic to) the tensor of $F_\eps$ and the writer
monad, which by construction forces commutativity; $\mathsf{FreeT}_\eps(\mathsf{Writer}_R)$ is
instead (up to the usual caveats about monad transformers not always agreeing with the
corresponding algebraic sum/tensor) closer to a \emph{sum}-like combination in which the two
signatures interleave positionally rather than being identified up to commutation --- see
Rivas and Jaskelioff's account of composing monoid-shaped notions of computation for the
general phenomenon.

\paragraph{Soundness and adequacy.} The original Theorems 5.3--5.6 are proved by induction on
the operational semantics, using at each step only: (i) the generic monad interface of
$S_\eps$; (ii) the fact that $R_\eps$ is an action $\eps$-algebra; and (iii) the specific
leaf-shape reasoning inside \texttt{loss}, $(\then)$, and \texttt{reset}, which we have redone
above against $\widehat W_\eps$. Since the operational semantics of $\lC$ is untouched by this
note, we expect analogues of Theorems 5.3--5.6 to hold for $\widehat S_\eps$ with $\Rm\eps$ and
$\collect/\collect_X/\censor$ in place of $R_\eps$ and the leaf-pattern-matching they replace;
the case analyses of Lemma~5.1 and~5.2 go through unchanged since they only concern values and
terminal expressions, which never touch the loss bookkeeping at all. The one place a genuinely
new argument would be needed is wherever the original soundness proof implicitly uses
commutativity of loss with operation calls to reorder a derivation step -- we have not located
such a step in the small-step soundness proof sketch, but a full re-verification of
Theorem~5.3 against $\widehat S_\eps$ is left as future work, alongside a from-scratch proof
that $\widehat W_\eps$ and $\widehat S_\eps$ as defined here indeed satisfy the monad laws (we
have checked the unit laws and the algebra-homomorphism property of $f\dg{\widehat W_\eps}$ in
Section~3.2, but not associativity of bind in full).

As a form of executable evidence, \texttt{Correct.hs} implements exactly $\widehat S_\eps$ (as
argued constructively above, definition by definition), and runs the same example programs as
the original selection-monad implementation, giving operational -- if not yet formally
proved -- confidence that the construction is coherent.

\section{Porting Appendix B: Correctness Proofs}

Appendix B of the Full Version repeats the denotational semantics (its Sections B.1--B.3
duplicate the material we have already redone in Sections 2--5 above) and then, in Section
B.4, proves soundness and adequacy against the operational semantics. We now port Section B.4
lemma by lemma. The operational semantics of $\lC$ (its reduction rules, frames, and the
termination theorem, Appendix A) is untouched by the swap, so every statement below is
word-for-word the original's statement with $S_\eps,W_\eps,R_\eps,F_\eps,\varphi,\delta_\eps$
replaced by $\widehat S_\eps,\widehat W_\eps,\Rm\eps,$ (no separate plain free monad),
$\widehat\varphi,\collect$; only the \emph{proofs} of the lemmas that manipulate the leaf shape
of $W_\eps$ need new arguments, using $\tell,\censor,\collect,\collectX{(-)}$ from
Section~3.3--3.4. As in the original, we restrict attention to the core calculus of Sections
2--5 (variables, constants, tuples, projection, abstraction, application, operations,
\texttt{loss}, $(\then)$, the local construct, \texttt{reset}, and handling); extending to sum
types, \texttt{nat}, and \texttt{list} is routine and, as in the original text, entirely
orthogonal to the swap (every one of those extra clauses is proved exactly as in the original,
verbatim, since none of them ever inspects the shape of $W_\eps$).

\subsection{Substitution and value semantics}

\begin{lemma}[Substitution]
Suppose $\Gamma\vdash v:\sigma$ and $\Gamma,x{:}\sigma\vdash e:\tau\,!\,\eps$. Then
$\Gamma\vdash e[v/x]:\tau\,!\,\eps$.
\end{lemma}

This is a syntactic fact about typing and is untouched by the swap. As before, we record the
value semantics $\Vsem{v}:\Sm\Gamma\to\Sm\sigma$ for $\Gamma\vdash v:\sigma$, needed to state
soundness and adequacy:
\[
  \Vsem{x}(\rho) = \rho(x), \qquad
  \Vsem{c}(\rho) = [\![c]\!], \qquad
  \Vsem{(v_1,\dots,v_n)}(\rho) = (\Vsem{v_1}(\rho),\dots,\Vsem{v_n}(\rho)), \qquad
  \Vsem{\lambda^{\eps_1}x{:}\sigma_1.\,e}(\rho) = \lambda a\in\Ssem{\sigma_1}.\ \Ssem{e}(\rho[a/x]).
\]

\subsection{Values denote units (Lemma B.2)}

\begin{lemma}
For any value $\Gamma\vdash v:\sigma\,!\,\eps$: $\Ssem{v}(\rho) = \eta^{\widehat S_\eps}(\Vsem{v}(\rho))$.
\end{lemma}
\begin{proof}
By structural induction on $v$; every case is generic in the monad and transfers with $\eta^{\widehat S_\eps}$
in place of $\eta^{S_\eps}$. E.g.\ for $v=x$: $\Ssem{x}(\rho) = \eta^{\widehat S_\eps}(\rho(x)) =
\eta^{\widehat S_\eps}(\Vsem{x}(\rho))$; for $v=(v_1,\dots,v_n)$, using the induction hypothesis
at each $v_i$ and that $\eta^{\widehat S_\eps}$ is a unit for $\mathsf{let}_{\widehat S_\eps}$,
\[
  \Ssem{(v_1,\dots,v_n)}(\rho) = \mathsf{let}_{\widehat S_\eps}\ a_1\in\Ssem{\sigma_1}\ \mathsf{be}\
     \eta^{\widehat S_\eps}(\Vsem{v_1}(\rho))\ \mathsf{in}\ \cdots\ \mathsf{in}\ \eta^{\widehat S_\eps}((a_1,\dots,a_n))
   = \eta^{\widehat S_\eps}\big((\Vsem{v_1}(\rho),\dots,\Vsem{v_n}(\rho))\big);
\]
and for $v=\lambda^{\eps_1}x{:}\sigma_1.e$, directly $\Ssem{v}(\rho) = \eta^{\widehat
S_\eps}(\lambda a.\Ssem{e}(\rho[a/x])) = \eta^{\widehat S_\eps}(\Vsem{v}(\rho))$. None of these
steps ever inspects $\widehat W_\eps$.
\end{proof}

\subsection{The action law simplifies (Lemma B.3)}

Fix a handler $h$, environment $\rho$, and $\gamma\in\Sm{\sigma'}\to\Rm\eps$ as in Section~5,
and let $s:\Sm\sigma\to A$ ($A = \widehat W_\eps(\Sm{\sigma'})^{\Sm{\mathit{par}}}$) be the
return-clause map from Section~5. Write $\tell(r)^P$ for $\tell(r)$ applied pointwise to a
function into $A$.

\begin{lemma}
$s\dg{\widehat W_{\eps\ell}}\circ\tell(r) = \tell(r)^P\circ s\dg{\widehat W_{\eps\ell}}$.
\end{lemma}
\begin{proof}
The original's proof of the analogous fact needs the extension-uniqueness argument (both sides
are shown to be homomorphic extensions of the same map on generators) precisely because there
$r\cdot(-)$ acts on \emph{every} leaf of $F_\eps(R\times X)$ simultaneously, and one must invoke
the commutativity of the action with $\psi$ to know that the extended map is still a
homomorphism. Here $\tell(r)$ only ever touches the \emph{root} loss slot -- $\tell(r)(r_0,n) =
(r+r_0,n)$ -- so the fact is immediate from the recursive definition of $f\dg{\widehat
W_{\eps\ell}}$ itself, with no commutation needed at all:
\[
  s\dg{\widehat W_{\eps\ell}}(\tell(r)(r_0,n)) = s\dg{\widehat W_{\eps\ell}}(r+r_0,n)
     = (r+r_0)\cdot(\cdots) = r\cdot\big(r_0\cdot(\cdots)\big) = r\cdot s\dg{\widehat W_{\eps\ell}}(r_0,n),
\]
using only additivity of the action on $A$ (the case split on whether $n$ is a leaf or an
internal node is the same on both sides of ``$(\cdots)$'' and immaterial to the argument). This
is one of the two places (with the return clause of Section~5) where the swap strictly
simplifies a piece of the original development.
\end{proof}

\subsection{Pairing and application (Lemma B.4)}

\begin{lemma}\
\begin{enumerate}[label=(\arabic*)]
\item For $\Gamma\vdash v:\sigma$ and $\Gamma\vdash e:\tau\,!\,\eps$:
  $\Ssem{(v,e)}(\rho) = \mathsf{let}_{\widehat S_\eps}\ b\in\Ssem\tau\ \mathsf{be}\ \Ssem{e}(\rho)\ \mathsf{in}\ \eta^{\widehat S_\eps}((\Vsem{v}(\rho),b))$.
\item For $\Gamma\vdash v:\sigma\to\tau\,!\,\eps$ and $\Gamma\vdash e:\sigma\,!\,\eps$:
  $\Ssem{v\,e}(\rho) = \mathsf{let}_{\widehat S_\eps}\ a\in\Ssem\sigma\ \mathsf{be}\ \Ssem{e}(\rho)\ \mathsf{in}\ \Vsem{v}(\rho)(a)$.
\end{enumerate}
\end{lemma}
\begin{proof}
Both are the generic clauses for tuples and application, using Lemma~5.1's hat-analogue (just
proved) to replace $\Ssem{v_i}(\rho)$ by $\eta^{\widehat S_\eps}(\Vsem{v_i}(\rho))$ and the unit
law for $\mathsf{let}_{\widehat S_\eps}$; neither uses $\widehat W_\eps$'s internal shape.
\end{proof}

\subsection{Loss continuations and $(\then)$ (Lemma B.5)}

This is the first place the internal shape of $\widehat W_\eps$ genuinely enters. Recall from
Section~4 the \emph{general} definition of the loss-function semantics, for any
$\Gamma,x{:}\sigma\vdash e:\mathtt{loss}\,!\,\eps$:
$\Lsem{\lambda^\eps x{:}\sigma.\,e}(\rho)(a) = \collapse(\Ssem{e}(\rho[a/x])(\gamma_0))$, where
$\gamma_0 = \lambda r{\in}R.\,\eta^{\widehat W_\eps}(())$. The lemma below is not a new
definition but a \emph{derived compositional law} for the special case where the body is
itself a $(\then)$-chain, exactly mirroring the role of the original Lemma B.5.

\begin{lemma}\
\begin{enumerate}[label=(\arabic*)]
\item For $\Gamma\vdash e:\sigma\,!\,\eps$ and $\Gamma\vdash g:\sigma\to\mathtt{loss}\,!\,\eps_1$
  with $\eps_1\subseteq\eps$, and for every $\gamma_1$:
  \[
    \Ssem{e\then g}(\rho)(\gamma_1) = \collect\big(\Rx\eps(\Ssem{e}(\rho)\mid\Lsem{g}(\rho))\big).
  \]
\item For $\Gamma\vdash g:\sigma\to\mathtt{loss}\,!\,\eps$:
  $\collect\circ\Lsem{g}(\rho) = \lambda a\in\Ssem\sigma.\ \Vsem{g}(\rho)\,a\,\gamma_1$ for every $\gamma_1$.
\item For $\Gamma,x{:}\sigma\vdash e:\tau\,!\,\eps$ and $\Gamma,x{:}\sigma\vdash g:\tau\to\mathtt{loss}\,!\,\eps_1$
  with $\eps_1\subseteq\eps$:
  \[
    \Lsem{\lambda^\eps x{:}\sigma.\,e\then g}(\rho) = \lambda a\in\Ssem\sigma.\ \Rx\eps\big(\Ssem{e}(\rho[a/x])\mid\Lsem{g}(\rho[a/x])\big).
  \]
\end{enumerate}
\end{lemma}
\begin{proof}
Parts (1) and (2) by mutual induction on $g$.

\emph{Part (1).} Unfolding the Section~4 definition of $(\then)$ (which manifestly does not
mention its own argument $\gamma_1$, giving the claimed constancy for free) and applying
$\collectX{\Ssem\sigma}$ to expose the pair of value and accumulated loss at the leaves of
$\Ssem{e}(\rho)(\Lsem{g}(\rho))$:
\[
  \Ssem{e\then g}(\rho)(\gamma_1) =
  \mathsf{let}_{\widehat W_\eps}\ (a,r_1)\in\Ssem\sigma\times R\ \mathsf{be}\ \collectX{\Ssem\sigma}\big(\Ssem{e}(\rho)(\Lsem{g}(\rho))\big)\ \mathsf{in}\
  \mathsf{let}_{\widehat W_\eps}\ (r_3,r_2)\in R\times R\ \mathsf{be}\ \collectX{R}\big(\Vsem{g}(\rho)\,a\,\gamma_0\big)\ \mathsf{in}\ \tell(r_2)\big(\eta^{\widehat W_\eps}(r_1+r_3)\big).
\]
Using the induction hypothesis (2) to replace $\Vsem{g}(\rho)\,a\,\gamma_0$ by
$\collect(\Lsem{g}(\rho)(a))$, and unfolding $\collect = \Wm\eps(\pi_2)\circ\collectX 1$ so that
$\collectX{R}(\collect(u)) = \collectX R(\Wm\eps(\pi_2)(\collectX 1(u)))$ has the same total
accumulated loss as $\collectX 1(u)$ paired with its (now doubly-recorded) value, one checks by
the defining clauses of $\collectX{(-)}$, $\bump$, and $\tell$ that the right-hand side above
collapses to exactly $\gamma^{\dagger\Wm\eps}(\Ssem{e}(\rho)(\gamma))$ for $\gamma := \Lsem
g(\rho)$, i.e.\ to $\Rx\eps(\Ssem e(\rho)\mid\Lsem g(\rho)) \in \Rm\eps$, both sides simply
re-adding the same two loss contributions $r_1$ and (via $g$) $r_3$ at the same point and
zeroing the same intermediate bookkeeping; applying the outer $\collect$ then promotes that
$\Rm\eps$ value to the $\Wm\eps(R)$-typed result, matching the left-hand side exactly. (This is
the hat-analogue of the original's five-line chain via $\delta_\eps$; here every step is a
direct unfolding of $\collect,\collectX{(-)},\tell$ rather than a bind-fusion law for
$F_\eps(R\times-)$'s $\delta_\eps$, but the computation performed is the same one: total up the
loss along $e$'s evaluation, total up the loss along $g(a)$'s evaluation against the zero
continuation, add, and record.)

\emph{Part (2).} If $g=\lambda^\eps x{:}\sigma.0$ both sides are $\collect(\eta^{\widehat
W_\eps}(0)) = \collect(0,\iota_1(()))=(0,\iota_1(0))$ against $\Vsem g(\rho)\,a\,\gamma_1 =
(0,\iota_1(0))$ directly. If $g=\lambda^\eps x{:}\sigma.\,e_2\then g_1$, then $\Vsem
g(\rho)\,a\,\gamma_1 = \Ssem{e_2\then g_1}(\rho[a/x])(\gamma_1) = \collect\big(\Rx\eps(\Ssem{e_2}(\rho[a/x])\mid\Lsem{g_1}(\rho[a/x]))\big)$
by part (1), which is exactly $\collect(\Lsem g(\rho)(a))$ by part (3) applied at $e:=e_2,g:=g_1$
-- so it remains to note part (3) does not depend on part (2) at the same $g$, only at the
strictly shorter $g_1$, closing the induction.

\emph{Part (3).} Immediate from parts (1)-(2): unfolding the base definition of $\mathcal L$ at
$e' := e\then g$ gives $\collapse(\Ssem{e\then g}(\rho[a/x])(\gamma_0))$; by part (1) this is
$\collapse(\collect(\Rx\eps(\Ssem e(\rho[a/x])\mid\Lsem g(\rho[a/x]))))$, and one checks directly
from the clauses of $\collapse$ and $\collect$ that $\collapse\circ\collect = \mathrm{id}$ on
$\Rm\eps$ (collect installs the total loss as a fresh $R$-valued leaf with everything else
zeroed; collapse discards the (now-zero) root loss and keeps that leaf value, recovering exactly
the argument), giving $\Rx\eps(\Ssem e(\rho[a/x])\mid\Lsem g(\rho[a/x]))$ as required.
\end{proof}

\subsection{The context lemma (Lemma B.6)}

\begin{lemma}
For $e:\sigma\,!\,\eps$ and $F[e]:\tau\,!\,\eps$ where $F$ is a regular frame:
\[
  \Ssem{F[e]}(\rho) = \mathsf{let}_{\widehat S_\eps}\ a\in\Ssem\sigma\ \mathsf{be}\ \Ssem{e}(\rho)\ \mathsf{in}\ \Ssem{F[x]}(\rho[x/a]).
\]
\end{lemma}
\begin{proof}
By cases on the form of $F$, as in the original; we record the same illustrative examples. For
$F=f(\square)$: both sides unfold to $\eta^{\widehat S_\eps}([\![f]\!](a))$ under the binding
$a:=\Ssem e(\rho)$, using only the unit law of $\mathsf{let}_{\widehat S_\eps}$. For $F=\square.i$:
$\Ssem{e.i}(\rho) = \widehat S_\eps(\pi_i)(\Ssem e(\rho)) = \mathsf{let}_{\widehat S_\eps}\ a\in(\sigma_1,\dots,\sigma_n)\
\mathsf{be}\ \Ssem e(\rho)\ \mathsf{in}\ \pi_i(a)$, again purely a fact about the abstract
$\widehat S_\eps$ interface. No case here touches $\widehat W_\eps$'s internal structure, so
every case transfers verbatim with $\widehat S_\eps$ in place of $S_\eps$.
\end{proof}

\subsection{Threading loss continuations through frames (Lemma B.7)}

\begin{lemma}
For $e:\sigma\,!\,\eps$, $F[e]:\tau\,!\,\eps$, and $g:\tau\to\mathtt{loss}\,!\,\eps$:
\[
  \Ssem{F[e]}\Lsem g = \mathsf{let}_{\widehat W_\eps}\ a\in\Ssem\sigma\ \mathsf{be}\ \Ssem{e}\Lsem{\lambda^\eps x{:}\sigma.\,F[x]\then g}\ \mathsf{in}\ \Ssem{F[x]}(x{\mapsto}a)\Lsem{g}.
\]
\end{lemma}
\begin{proof}
Exactly the original's three-line argument, unchanged in form:
\[
  \Ssem{F[e]}\Lsem g
  \overset{\text{Lem.~5.6(hat)}}{=} \big(\mathsf{let}_{\widehat S_\eps}\ a\in\Ssem\sigma\ \mathsf{be}\ \Ssem e\ \mathsf{in}\ \Ssem{F[x]}(x{\mapsto}a)\big)\Lsem g
  \overset{\text{Eq.\ for } f\dg{\widehat S_\eps}}{=} \mathsf{let}_{\widehat W_\eps}\ a\in\Ssem\sigma\ \mathsf{be}\ \Ssem e\big(\lambda a.\,\Rx\eps(\Ssem{F[x]}(x{\mapsto}a)\mid\Lsem g)\big)\ \mathsf{in}\ \Ssem{F[x]}(x{\mapsto}a)\Lsem g,
\]
and the inner continuation $\lambda a.\,\Rx\eps(\Ssem{F[x]}(x{\mapsto}a)\mid\Lsem g)$ is exactly
$\Lsem{\lambda^\eps x{:}\sigma.\,F[x]\then g}$ by Lemma~7.5(3) (the hat-analogue of Lemma~B.5),
giving the claim. Since the Kleisli-extension formula for $\widehat S_\eps$ is syntactically
identical to the original's equation~(6) (Section~3.4), this proof is a verbatim transcription.
\end{proof}

\subsection{Operations under contexts (Lemma B.8)}

\begin{lemma}
For $K[\op(v)]:\sigma\,!\,\eps$ with $\op\notin\mathrm{heff}(K)$ and $\op:\mathit{out}\xrightarrow{\ell}\mathit{in}$:
\[
  \Ssem{K[\op(v)]}(\rho)(\gamma) = \widehat\varphi_{\ell,\op,\eps(\ell)}\big(\Vsem v(\rho),\ \lambda a\in\Ssem{\mathit{in}}.\ \Ssem{K[x]}(\rho[a/x])(\gamma)\big).
\]
\end{lemma}
\begin{proof}
By induction on the size of $K$. The base case $K=\square$ is the $\op(e)$ clause of Figure~9
(hat-version) applied to a value: $\Ssem{\op(v)}(\rho)(\gamma) = \widehat\varphi_{\ell,\op,\eps(\ell)}(\Vsem
v(\rho),(\eta^{\widehat W_\eps})^{\Ssem{\mathit{in}}})$, and $(\eta^{\widehat
W_\eps})^{\Ssem{in}} = \lambda a.\,(0,\iota_1(a)) = \lambda a.\,\Ssem{x}(\rho[a/x])(\gamma)$ for
any $\gamma$ (since $\Ssem x(\rho')(\gamma)=\eta^{\widehat W_\eps}(\rho'(x))$ regardless of
$\gamma$), matching the claim. For the inductive cases:
\begin{itemize}
\item $K=F[K_1]$ (regular frame): purely a fact about the abstract $\widehat S_\eps$-algebra
  structure of $\widehat\varphi^X$ combined with Lemma~7.6 (hat-Lemma~B.6), unchanged in form
  from the original: the induction hypothesis lets us pull $F[\square]$ through
  $\widehat\varphi^{\Ssem\tau}_{\ell,\op,\eps(\ell)}$, using that $\widehat\varphi^X_{\ell,\op,i}(o,f)(\gamma)
  = \widehat\varphi_{\ell,\op,i}(o,\lambda x.f(x)(\gamma))$ is exactly the naturality needed.
\item $K=\mathtt{with}\ h\ \mathtt{from}\ v\ \mathtt{handle}\ K_1$, handling $\ell'\ne\ell$:
  writing $T[\gamma]=\lambda a.\,\gamma\dg{\widehat W_\eps}(s(a)(\Vsem v(\rho)))$ for the
  continuation constructed in Section~5's handler formula, we have by that formula and the
  induction hypothesis
  \[
    \Ssem{\mathtt{with}\ h\ \mathtt{from}\ v\ \mathtt{handle}\ K_1[\op(v)]}(\rho)(\gamma)
       = s\dg{\widehat W_{\eps\ell'}}\Big(\widehat\varphi_{\ell,\op,\eps(\ell)}\big(\Vsem v(\rho),\lambda a.\,\Ssem{K_1[x]}(\rho[a/x])(T[\gamma])\big)\Big)(\Vsem v(\rho)),
  \]
  and since $\ell\ne\ell'$ (the operation is not the one $h$ handles), $\widehat\varphi_{\ell,\op,\eps(\ell)}$
  is exactly the pass-through case of $\widehat\psi$ from Section~5, so $s\dg{\widehat
  W_{\eps\ell'}}$ applied to it unfolds, by the defining recursion of $\dg{\widehat
  W_{\eps\ell'}}$, to $\widehat\varphi_{\ell,\op,\eps(\ell)}$ again (now one level further out),
  giving exactly $\widehat\varphi_{\ell,\op,\eps(\ell)}(\Vsem v(\rho), \lambda a.\, \Ssem{\mathtt{with}\ h\ \mathtt{from}\ v\ \mathtt{handle}\ K_1[x]}(\rho[a/x])(\gamma))$
  as required.
\item $K=\loc{K_1}g^{\eps_1}$: the local-construct clause $\Ssem{\loc e g^{\eps_1}}(\rho)=\lambda\gamma.\Ssem e(\rho)(\Lsem g(\rho))$
  simply substitutes a fixed continuation and commutes with $\widehat\varphi$ trivially (it does
  not touch $\widehat W_\eps$'s shape at all), so this case is immediate from the induction
  hypothesis.
\item $K=K_1\then\lambda^{\eps_1}x{:}\tau.\,e_1$: unfolding Section~4's $(\then)$-clause and the
  induction hypothesis on $K_1$, the node $\widehat\varphi_{\ell,\op,\eps(\ell)}(o,\kappa)$
  produced by $K_1$ carries recorded loss $0$ at its own root (by definition of $\widehat\varphi$),
  so $\collectX{\Ssem\tau}$ and $\tell$ both simply recurse through it into its children
  unchanged, letting the whole $(\then)$-construction be pulled inside the
  $\lambda a\in\Ssem{\mathit{in}}.\,(-)$ exactly as in the original's calculation (there, using
  that $F_\eps(R\times-)$'s bind passes through internal nodes verbatim; here, using that
  $\collectX{(-)}$'s and $\tell$'s Free-node clauses are literally the identity-on-shape
  recursions given in Section~3.3--3.4).
\item $K=\mathtt{reset}\ K_1$: by the $\censor$-based clause of Section~4, and because
  $\censor$'s own Free-node clause is $\censor(r_0,\iota_2((\ell,\op,i),(o,\kappa))) =
  (0,\iota_2((\ell,\op,i),(o,\lambda a.\censor(\kappa(a)))))$ -- i.e.\ $\censor$ commutes with
  $\widehat\varphi_{\ell,\op,i}$ by its very definition, not merely up to a lemma -- the case
  is immediate: $\censor(\widehat\varphi_{\ell,\op,\eps(\ell)}(\Vsem v(\rho),\kappa)) =
  \widehat\varphi_{\ell,\op,\eps(\ell)}(\Vsem v(\rho),\lambda a.\censor(\kappa(a)))$, matching
  $\Ssem{\mathtt{reset}\ K_1[x]}(\rho[a/x])(\gamma)$ pointwise.
\end{itemize}
\end{proof}

\subsection{Small-step soundness (Theorem B.9)}

\begin{theorem}[Small-step soundness]
Suppose $e:\sigma\,!\,\eps$ and $g:\sigma\to\mathtt{loss}\,!\,\eps_1$ with $\eps_1\subseteq\eps$. Then
\[
  g\vdash_\eps e\xrightarrow{r} e' \implies \Ssem{e}\Lsem g = r\cdot\big(\Ssem{e'}\Lsem g\big).
\]
\end{theorem}
\begin{proof}
By cases on the transition, using Lemma~7.2 (hat-Lemma~B.2) without comment. Cases (1)-(8) of
the original (basic function application, projection, \texttt{cases}/\texttt{iter}/\texttt{fold}
reduction, $\beta$-reduction) never mention $F_\eps$, $W_\eps$, or $R_\eps$ at all -- they only
manipulate $\eta^{S_\eps}$, $\mathsf{let}_{S_\eps}$, and value semantics -- so they transfer
completely unchanged with $\widehat S_\eps,\eta^{\widehat S_\eps}$ in place of $S_\eps,\eta^{S_\eps}$;
e.g.\ case (1), $f(v)\xrightarrow{0}v'$ where $f(v)\to v'$:
\[
  \Ssem{f(v)}\Lsem g = \mathsf{let}_{\widehat S_\eps}\ a\in\Ssem\sigma\ \mathsf{be}\ \eta^{\widehat S_\eps}(\Vsem v)\ \mathsf{in}\
    \eta^{\widehat S_\eps}([\![f]\!](a)) = \eta^{\widehat S_\eps}(\Vsem{v'}) = \Ssem{v'}\Lsem g = 0\cdot\Ssem{v'}\Lsem g.
\]
We give the remaining, swap-relevant cases in full.

\emph{(9) $g\vdash_\eps \mathtt{loss}(r)\xrightarrow{r} ()$.} For any $\gamma$:
\[
  \Ssem{\mathtt{loss}(r)}(\gamma) = \mathsf{let}_{\widehat W_\eps}\ a\in R\ \mathsf{be}\ \eta^{\widehat W_\eps}(r)\ \mathsf{in}\ \tell(a)(\eta^{\widehat W_\eps}(()))
     = \tell(r)(\eta^{\widehat W_\eps}(())) = (r,\iota_1(())) = r\cdot(0,\iota_1(())) = r\cdot\Ssem{()}(\gamma).
\]

\emph{(13) $g\vdash_\eps v\then\lambda^{\eps_1}x{:}\sigma.e\xrightarrow{0}\loc{e[v/x]}{\lambda^{\eps\eps_1}x{:}\sigma.0}$.}
Unfolding Section~4's clause with $\Ssem v(\rho)=\eta^{\widehat W_\eps}(\Vsem v)$:
\[
  \Ssem{v\then\lambda x.e}(\gamma) = \mathsf{let}_{\widehat W_\eps}(a,r_1)\in\sigma\times R\ \mathsf{be}\ \collectX\sigma(0,\iota_1(\Vsem v))\ \mathsf{in}\ \cdots
     = \mathsf{let}_{\widehat W_{\eps_1}}(r_3,r_2)\in R\times R\ \mathsf{be}\ \collectX{R}\big(\Ssem e(x{\mapsto}\Vsem v)(\gamma_0)\big)\ \mathsf{in}\ \tell(r_2)(\eta(0+r_3)),
\]
since $\collectX\sigma(0,\iota_1(\Vsem v)) = (0,\iota_1((\Vsem v,0)))$ forces $a=\Vsem v$,
$r_1=0$. This is precisely $\collapse(\Ssem{e[v/x]}(\gamma_0))$ up to relabelling $r_2,r_3$,
which is $\Lsem{\lambda^{\eps\eps_1}x{:}\sigma.0}$'s value at $()$ composed with the general
$\mathcal L$-definition applied to $e[v/x]$ -- i.e.\ exactly $\Ssem{\loc{e[v/x]}{\lambda x.0}}(\gamma)$
by the local-construct clause, matching the target with loss $0$.

\emph{(14) $g\vdash_\eps \mathtt{reset}\ v\xrightarrow{0}v$.}
$\Ssem{\mathtt{reset}\ v}\Lsem g = \censor(\Ssem v\Lsem g) = \censor(0,\iota_1(\Vsem v)) = (0,\iota_1(\Vsem v)) = \Ssem v\Lsem g = 0\cdot\Ssem v\Lsem g$,
using that $\censor$ is already the identity on a $\iota_1$-leaf carrying recorded loss $0$.

\emph{(15) Frame rule} ($\lambda^\eps x{:}\tau.(F[x]\then g)\vdash_\eps e_1\xrightarrow{r}e_1'
\implies g\vdash_\eps F[e_1]\xrightarrow{r}F[e_1']$): by the induction hypothesis and Lemma~7.7
(hat-Lemma~B.7) applied twice, exactly as in the original:
\[
  \Ssem{F[e_1]}\Lsem g \overset{\text{Lem 7.7}}{=} \mathsf{let}_{\widehat W_\eps}\,a\,\mathsf{be}\,\Ssem{e_1}\Lsem{\lambda x.F[x]\then g}\,\mathsf{in}\,\Ssem{F[x]}(a)\Lsem g
     \overset{\text{IH}}{=} \mathsf{let}_{\widehat W_\eps}\,a\,\mathsf{be}\,r\cdot\Ssem{e_1'}\Lsem{\lambda x.F[x]\then g}\,\mathsf{in}\,\Ssem{F[x]}(a)\Lsem g
     = r\cdot\Ssem{F[e_1']}\Lsem g,
\]
the last step by $\tell(r)$ (i.e.\ $r\cdot(-)$) commuting with $\mathsf{let}_{\widehat W_\eps}$'s
leaf clause, which is immediate from the bind formula of Section~3.2.

\emph{(16) $\mathtt{with}\ h\ \mathtt{from}\ v\ \mathtt{handle}\ e_2\to \mathtt{with}\ h\ \mathtt{from}\ v\ \mathtt{handle}\ e_2'$
(step inside a handled body).} Let $v_r=\lambda^\eps z{:}(\mathit{par},\sigma_1).e_r$ be $h$'s
return clause. By the induction hypothesis, $\Ssem{e_2}\Lsem{\lambda^\eps x{:}\sigma_1.\,v_r(v,x)\then g} = r\cdot\Ssem{e_2'}\Lsem{\lambda^\eps x{:}\sigma_1.\,v_r(v,x)\then g}$.
As in the original, one checks $\Lsem{\lambda^\eps x{:}\sigma_1.v_r(v,x)\then g} = \lambda a.\,\Rx\eps(\Ssem{e_r}[(\Vsem v,a)/z]\mid\Lsem g)$
by Lemma~7.5(3), so both $\Ssem{\mathtt{with}\ h\ \mathtt{from}\ v\ \mathtt{handle}\ e_2}\Lsem g$
and its primed counterpart equal $s\dg{\widehat W_{\eps\ell}}(\Ssem{e_2}(\cdots))(\Vsem v)$ and
$s\dg{\widehat W_{\eps\ell}}(\Ssem{e_2'}(\cdots))(\Vsem v)$ respectively, by Section~5's handler
formula. Applying $s\dg{\widehat W_{\eps\ell}}$ (a function, hence respecting the induction
hypothesis's equation) and then Lemma~7.3 (hat-Lemma~B.3) to pull $r\cdot(-)$ back out:
\[
  s\dg{\widehat W_{\eps\ell}}\big(\Ssem{e_2}(\cdots)\big)(\Vsem v) = s\dg{\widehat W_{\eps\ell}}\big(r\cdot\Ssem{e_2'}(\cdots)\big)(\Vsem v)
  \overset{\text{Lem.\ 7.3}}{=} r\cdot\Big(s\dg{\widehat W_{\eps\ell}}\big(\Ssem{e_2'}(\cdots)\big)(\Vsem v)\Big),
\]
as required. (Recall from Section~7.3 that Lemma~7.3's proof no longer needs the original's
uniqueness-of-extension argument, so this case is, if anything, slightly more direct than the
original.)

\emph{(10) Operation call under a handler} ($g\vdash_\eps \mathtt{with}\ h\ \mathtt{from}\ v_1\ \mathtt{handle}\ K[\op(v_2)]\xrightarrow{0}v_o(v_1,v_2,f_l,f_k)$,
where $\op\notin\mathrm{heff}(K)$, $\op\mapsto v_o\in h$, and, crucially,
\[
  f_k = \lambda^\eps(p,y){:}(\mathit{par},\mathit{in}).\ \loc{\mathtt{with}\ h\ \mathtt{from}\ p\ \mathtt{handle}\ K[y]}{g}^\eps,
  \qquad
  f_l = \lambda^\eps(p,y){:}(\mathit{par},\mathit{in}).\ (\mathtt{with}\ h\ \mathtt{from}\ p\ \mathtt{handle}\ K[y])\then g
\]
-- note $f_k$ is wrapped in the local construct $\loc{\cdot}{g}^\eps$; \emph{this wrapper is
essential, and Section~8 explains exactly why}). Using Lemma~7.8 (hat-Lemma~B.8) and writing
$T[\gamma]=\lambda a.\,\Rx\eps(\Ssem{e_r}[(\Vsem{v_1},a)/z]\mid\gamma)$:
\[
  \Ssem{\mathtt{with}\ h\ \mathtt{from}\ v_1\ \mathtt{handle}\ K[\op(v_2)]}\Lsem g
  = s\dg{\widehat W_{\eps\ell}}\Big(\widehat\varphi_{\ell,\op,i}\big(\Vsem{v_2},\lambda a.\,\Ssem{K[y]}(y{\mapsto}a)(T[\Lsem g])\big)\Big)(\Vsem{v_1})
  = \Ssem{e_o}\big(z{\mapsto}(\Vsem{v_1},\Vsem{v_2},l_1,k_1)\big)\Lsem g,
\]
where $k_1(p,a) = s\dg{\widehat W_{\eps\ell}}(\Ssem{K[y]}(y{\mapsto}a)(T[\Lsem g]))(p)$ and
$l_1(p,a) = \collect\big(\Lsem g\dg{\widehat W_\eps}(k_1(p,a))\big)$ (a bare, $\gamma$-independent
value of $\Sm\eps(R)$, since $\collect$ and $\Lsem g\dg{\widehat W_\eps}$ have already discharged
all dependence on any further continuation) -- and one checks, exactly as in the original,
\[
  \Ssem{f_k}(p,a)(\gamma_1) = \Ssem{\loc{\mathtt{with}\ h\ \mathtt{from}\ p\ \mathtt{handle}\ K[y]}{g}^\eps}(y{\mapsto}a)(\gamma_1)
     = \Ssem{\mathtt{with}\ h\ \mathtt{from}\ p\ \mathtt{handle}\ K[y]}(y{\mapsto}a)\Lsem g = k_1(p,a)(\gamma_1),
\]
\emph{for every $\gamma_1$}, since the local-construct clause $\Ssem{\loc e g^\eps}(\rho)=\lambda\gamma.\Ssem e(\rho)(\Lsem g(\rho))$
overrides whatever continuation $f_k$ is applied to with $\Lsem g$ -- this is exactly why
$S[|f_k|](p,a)$ can be identified with the single, $\gamma$-independent element $k_1(p,a)$ of
$A=\widehat W_\eps(\Sm{\sigma'})^{\Sm{\mathit{par}}}$ that the handler algebra of Section~5
requires. Similarly $\Ssem{f_l}(p,a)(\gamma_1) = \collect(\Rx\eps(\Ssem{\mathtt{with}\ h\
\mathtt{from}\ p\ \mathtt{handle}\ K[y]}(y{\mapsto}a)\mid\Lsem g)) = \collect(\Lsem g\dg{\widehat
W_\eps}(k_1(p,a))) = l_1(p,a)$, by Lemma~7.5(1) followed by the identity just established.

\emph{(11) $g\vdash_\eps \mathtt{with}\ h\ \mathtt{from}\ v_1\ \mathtt{handle}\ v_2\xrightarrow{0}v_r(v_1,v_2)$.}
Using Lemma~7.2 to replace $\Ssem{v_2}$ by $\eta^{\widehat W_\eps}(\Vsem{v_2})$ inside the
handler formula of Section~5, and that $\tell(0)=\mathrm{id}$:
\[
  \Ssem{\mathtt{with}\ h\ \mathtt{from}\ v_1\ \mathtt{handle}\ v_2}(\gamma)
   = s\dg{\widehat W_{\eps\ell}}\big(\eta^{\widehat W_{\eps\ell}}(\Vsem{v_2})\big)(\Vsem{v_1})
   = 0\cdot s(\Vsem{v_2})(\Vsem{v_1}) = \Ssem{e_r}(z{\mapsto}(\Vsem{v_1},\Vsem{v_2}))(\gamma) = \Ssem{v_r(v_1,v_2)}(\gamma),
\]
using the return clause $s(a)=\lambda p.\Ssem{e_r}(\rho[(p,a)/z])(\gamma)$ of Section~5 directly
(with no $r\cdot(-)$ to carry, as noted in Section~5).

\emph{(17) $\loc{e}{g_1}^{\eps_1}$ frame rule.} By the induction hypothesis
$\Ssem e\Lsem{g_1} = r\cdot\Ssem{e'}\Lsem{g_1}$, and directly from the local-construct clause:
$\Ssem{\loc e{g_1}^{\eps_1}}\Lsem g = \Ssem e\Lsem{g_1} = r\cdot\Ssem{e'}\Lsem{g_1} = r\cdot\Ssem{\loc{e'}{g_1}^{\eps_1}}\Lsem g$.

\emph{(18) $g\vdash_\eps(e_1\then g_1)\xrightarrow{0}r_0+(e_1'\then g_1)$, where
$g_1\vdash_\eps e_1\xrightarrow{r_0}e_1'$.} By the induction hypothesis
$\Ssem{e_1}\Lsem{g_1}=r_0\cdot\Ssem{e_1'}\Lsem{g_1}$; unfolding Section~4's $(\then)$-clause and
using that $\collectX\sigma$ commutes with $\tell(r_0)$ (equivalently, $r_0\cdot(-)$ and
$\collectX\sigma$ both only add to, respectively expose, the same accumulated total, so
$\collectX\sigma(\tell(r_0)(u)) = $ the leaf-wise sum of $\collectX\sigma(u)$ shifted by $r_0$,
by the $\bump$ identity $\bump_{r_0}$ recorded in Section~3.3--3.4):
\[
  \Ssem{e_1\then g_1}\Lsem g = \mathsf{let}_{\widehat W_\eps}(a,r_1)\,\mathsf{be}\,\collectX\sigma(r_0\cdot\Ssem{e_1'}\Lsem{g_1})\,\mathsf{in}\,(\cdots)
    = \mathsf{let}_{\widehat W_\eps}(a,r_1)\,\mathsf{be}\,\collectX\sigma(\Ssem{e_1'}\Lsem{g_1})\,\mathsf{in}\,\tell(r_2)(\eta(r_0+r_1+r_3))
    = \Ssem{r_0+(e_1'\then g_1)}\Lsem g,
\]
matching the operational semantics' convention that $r_0+(-)$ prepends $r_0$ to whatever loss
the remaining expression eventually records.

\emph{(19) $\mathtt{reset}$ frame rule.} By the induction hypothesis
$\Ssem{e_1}\Lsem g = r\cdot\Ssem{e_1'}\Lsem g$, and since $\censor(\tell(r)(u)) = \censor(u)$ for
any $r$ (censoring discards whatever the root loss is, whether or not it was just incremented):
\[
  \Ssem{\mathtt{reset}\ e_1}\Lsem g = \censor(\Ssem{e_1}\Lsem g) = \censor\big(r\cdot\Ssem{e_1'}\Lsem g\big) = \censor\big(\Ssem{e_1'}\Lsem g\big) = \Ssem{\mathtt{reset}\ e_1'}\Lsem g = 0\cdot\Ssem{\mathtt{reset}\ e_1'}\Lsem g.
\]
\end{proof}

\subsection{Evaluation soundness, adequacy, and giant-step adequacy}

The remaining results are, essentially verbatim, consequences of Theorem~7.9 (hat-Theorem~B.9)
by exactly the original's arguments, because -- as observed at the end of Section~5 -- the
\emph{observable, top-level} shape of $\Ssem e\Lsem g$ for a fully- or partially-evaluated $e$ is
literally identical in the two semantics: a pair $(r,v)$ when $e$ evaluates to a value (the root
loss slot of a $\iota_1$-leaf plays exactly the role of the original's leaf-pair), or
$\widehat\varphi_{\ell,\op,\eps(\ell)}(v,f)$ with $r$ folded into that node's own root-loss slot
via $\tell(r)$ when $e$ evaluates to a term stuck on $\op$. None of Theorems~B.10, B.11,
Corollary~B.12, or Theorem~B.13 inspect any \emph{deeper} structure of $\widehat W_\eps$ than
this top-level pair, so:

\begin{theorem}[Evaluation soundness]
For $e:\sigma\,!\,\eps$ and $g:\sigma\to\mathtt{loss}\,!\,\eps'$ with $\eps'\subseteq\eps$:
\[
  g\vdash_\eps e\xRightarrow{r}v \implies \Ssem e\Lsem g = (r,\Vsem v), \qquad
  g\vdash_\eps e\xRightarrow{r}K[\op(v)] \implies \Ssem e\Lsem g = \widehat\varphi_{\ell,\op,\eps(\ell)}\big(\Vsem v,\lambda a.\,r\cdot\Ssem{K[x]}(x{\mapsto}a)\Lsem g\big).
\]
\end{theorem}
\begin{proof}
By induction on the big-step derivation, exactly as the original: in the base case $e=v$, $r=0$,
Lemma~7.2 gives $\Ssem e\Lsem g = (0,\iota_1(\Vsem v))$; in the base case $e=K[\op(v)]$, $r=0$,
Lemma~7.8 gives the second equation directly with $\tell(0)=\mathrm{id}$; in both inductive
cases, the small-step case $g\vdash_\eps e\xrightarrow{r_1}e'$, $g\vdash_\eps
e'\xRightarrow{r_2}(\cdots)$, $r=r_1+r_2$ combines with Theorem~7.9 and the additivity of
$\tell$ ($\tell(r_1)\circ\tell(r_2) = \tell(r_1+r_2)$) to give the result.
\end{proof}

\begin{theorem}[Adequacy]
For $e:\sigma\,!\,\eps$ and $g:\sigma\to\mathtt{loss}\,!\,\eps'$ with $\eps'\subseteq\eps$:
\[
  \Ssem e\Lsem g = (r,a) \implies \exists v.\ g\vdash_\eps e\xRightarrow{r}v \wedge \Vsem v = a, \qquad
  \Ssem e\Lsem g = \widehat\varphi_{\ell,\op,\eps(\ell)}(a,f) \implies \exists K[\op(v)].\ g\vdash_\eps e\xRightarrow{r}K[\op(v)]\ \wedge\ a=\Vsem v\ \wedge\ f=\lambda b.\,r\cdot\Ssem{K[x]}(x{\mapsto}b)\Lsem g.
\]
\end{theorem}
\begin{proof}
Verbatim the original's argument from termination (Appendix~A, untouched by the swap) plus
Theorem~7.10: termination gives one of the two evaluation shapes; Theorem~7.10 rules out the
mismatching one (a $\iota_1$-leaf can never equal a $\widehat\varphi$-node, since $\widehat
W_\eps(X)$ is a disjoint sum) and identifies $r,a$ in the matching one.
\end{proof}

\begin{corollary}[First-order adequacy]
For first-order $\sigma$, $e:\sigma\,!\,\emptyset$, $g:\sigma\to\mathtt{loss}\,!\,\emptyset$, and
constants interpreted injectively: $\Ssem e\Lsem g = (r,\Vsem v) \iff g\vdash_\eps e\xRightarrow{r}v$.
\end{corollary}

Finally, giant-step adequacy transfers with the relation $\rhd$ between effect values $\mathrm{EV}$
and $\widehat W_\eps(\Ssem\sigma)$ redefined at the leaf level only: $(r,v)\rhd(s,\iota_1(a))
\iff r=s\wedge\Vsem v=a$, and $((\ell,\op),(v,f))\rhd(r,\iota_2((\ell',\op',n'),(a,g))) \iff
\ell=\ell'\wedge\op=\op'\wedge n'=\eps(\ell)\wedge\Vsem v=a\wedge\forall w.\,f(w)\rhd g(\Vsem w)$
(with $r$ left unconstrained on the operation-node clause, matching $\tell(r)$'s role of
recording the loss up to that node rather than the original's separate leaf-pairing).

\begin{theorem}[Giant-step adequacy]
For all $e:\sigma\,!\,\eps$: $\mathrm{Eval}(e) \rhd \Ssem e\Lsem g$.
\end{theorem}
\begin{proof}
By well-founded induction on $\mathrm{Eval}(e)$ exactly as the original, using
Theorem~7.10 in place of Theorem~B.10 in both cases.
\end{proof}

\section{An Error in the Original Appendix B.4}

We were asked to look for a mistake in the original proofs, and found one. Pages 49--51 of the
Full Version PDF (arXiv:2504.03890v1) contain, inside the proof of Theorem~B.9, \emph{two}
back-to-back derivations of case~(10) -- the operation-call-under-a-handler case -- for what is
presented as the same transition rule. The first derivation is complete and self-consistent; the
second is abandoned mid-argument at a line the authors themselves flag with the word ``BUT'',
and the text then moves directly on to case~(11) without resolving the discrepancy.

\paragraph{The first derivation (correct).} It uses
\[
  f_k = \lambda^\eps(p,y){:}(\mathit{par},\mathit{in}).\ \langle\mathtt{with}\ h\ \mathtt{from}\ p\ \mathtt{handle}\ K[y]\rangle g^\eps,
  \qquad
  f_l = \lambda^\eps(p,y){:}(\mathit{par},\mathit{in}).\ (\mathtt{with}\ h\ \mathtt{from}\ p\ \mathtt{handle}\ K[y])\blacktriangleright g,
\]
i.e.\ $f_k$ is wrapped in the local construct $\langle\cdot\rangle g^\eps$. This matches the
operational semantics rule as stated in the main body of the paper (the rule labelled with
$f_k,f_l$ defined this way, with the $\langle\cdot\rangle g^\eps$ wrapper on $f_k$). The
derivation closes cleanly with $S[|f_k|](p,a)\gamma_1 = k_1(p,a)\gamma_1$ and
$S[|f_l|](p,a)\gamma_1 = l_1(p,a)\gamma_1$ for every $\gamma_1$ -- exactly the fact we reprove
as our case~(10) in Section~7.9 above, and exactly because the $\langle\cdot\rangle g^\eps$
wrapper makes $S[|f_k|](p,a)$ \emph{constant} in its own continuation argument (Section~4's
local-construct clause always substitutes $L[|g|](\rho)$ regardless of what continuation
$f_k$ is applied to).

\paragraph{The second, abandoned derivation.} Immediately afterwards, the same transition is
redone with $f_k,f_l$ redefined \emph{without} the wrapper:
\[
  f_k = \lambda^\eps(p,y){:}(\mathit{par},\mathit{in}).\ \mathtt{with}\ h\ \mathtt{from}\ p\ \mathtt{handle}\ K[y],
  \qquad
  f_l = \lambda^\eps(p,y){:}(\mathit{par},\mathit{in}).\ f_k(p,y)\blacktriangleright g.
\]
Carrying through the same style of calculation, the authors derive
\[
  k_1(p,a) = \lambda\gamma.\ s\dg{F_{\eps\ell}}\big(S[|K[y]|](y{\mapsto}a)\,T(g)\big)\,p\,\gamma
\]
for the delimited continuation constructed inside the handler algebra (a function that, once
$g$ is fixed, is \emph{constant} in the further variable $\gamma$ -- it never uses it), flag the
line with ``BUT'', and then compute $S[|f_k|](p,a)$ directly from the semantics of the (now
unwrapped) expression $f_k$, obtaining
\[
  S[|f_k|](p,a) = \lambda\gamma.\ s\dg{F_{\eps\ell}}\big(S[|K[y]|](y{\mapsto}a)\big)(T(\gamma))\,p\,\gamma,
\]
which genuinely \emph{does} depend on $\gamma$ (through $T(\gamma)$, not the fixed $T(g)$) and
does not match the first expression. The manuscript stops right there, mid-derivation, and
begins case~(11).

\paragraph{Diagnosis.} This is not a transcription slip: we independently confirmed both
passages appear on consecutive pages of the source PDF exactly as described (a first, complete
proof using the wrapped $f_k$; then, with no transitional sentence, a second attempt using an
unwrapped $f_k$ that contradicts itself and stops). The mathematical content of the
contradiction is real and is precisely the point we relied on in our own case~(10)
(Section~7.9): \emph{without} the $\langle\cdot\rangle g^\eps$ wrapper, the semantics of the
reified delimited continuation $f_k$ genuinely depends on whatever loss continuation it is later
supplied with, so no single, $\gamma$-independent element $k_1(p,a)$ of the handler-algebra
carrier $A = W_\eps(\Sm{\sigma'})^{\Sm{\mathit{par}}}$ (equivalently $\widehat
W_\eps(\Sm{\sigma'})^{\Sm{\mathit{par}}}$ in our version) can represent it correctly. The
wrapper is not a stylistic choice; it is exactly what is needed for the compact,
$\gamma$-independent handler-algebra construction of Section~5 (original or ours) to be sound
at all for the delimited continuation captured at an operation call. We read the abandoned
second derivation as the authors' own record of discovering this necessity in the course of
writing the appendix -- valuable as a trace of \emph{why} the calculus's actual reduction rule
includes the wrapper, but it should not have been left in the ``Full Version'' as an
unresolved, self-contradicting proof fragment. Since our own handler semantics (Section~5) has
the same shape ($A=\widehat W_\eps(\Sm{\sigma'})^{\Sm{\mathit{par}}}$, built via a
$\gamma$-independent extension), the same diagnosis applies verbatim to our version: the swap
does not remove the need for the wrapper, it is simply orthogonal to it, since the wrapper is a
fact about the \emph{operational} semantics of $\lC$, which we have not touched.

\begin{thebibliography}{9}

\bibitem{plotkin-xie}
Gordon Plotkin and Ningning Xie.
\emph{Handling the Selection Monad (Full Version)}.
arXiv:2504.03890, 2025.

\bibitem{hyland}
Martin Hyland, Gordon Plotkin, and John Power.
\emph{Combining effects: Sum and tensor}.
Theoretical Computer Science, 357(1--3):70--99, 2006.

\bibitem{rivas}
Exequiel Rivas and Mauro Jaskelioff.
\emph{Notions of Computation as Monoids}.
Journal of Functional Programming, 27:e21, 2017.

\bibitem{pretnar}
Gordon Plotkin and Matija Pretnar.
\emph{Handling Algebraic Effects}.
Logical Methods in Computer Science, 9(4), 2013.

\end{thebibliography}

\end{document}
